\scanpage{92}
Si $k$ n'est pas de car. $0$, on ne peut en général aller en sens inverse, et $R(\mathbb{N},k)\xrightarrow{\Lambda}k[[t]]$ n'est \deleted{en général} pas injectif, [$k$ corps de car $>0$]

\medskip\centerline{\rule{2.6em}{0.3pt}}\medskip

On voit \deleted{que} par la formule de Lefschetz-Weil-Grothendieck que la connaissance de la fonction $L$ de $(X,F$ \deleted{\unsure{sur}}compactifiable sur $\mathbb{F}_p$, \inserted{à coeff. dans $F$,} \uline{équivaut} à celle de l'élément
\[
\underline{\chi}(F)=\sum(-1)^i\,\mathrm{cl}\,H^i_!(\overline{X},\overline{F})
\quad\text{de}\quad R(\mathbb{N},\mathbb{Q}_\ell)
\]
[car \unsure{essentiellt|essentiellement} cette fonction $L$ est l'image de $\underline{\chi}(F)$ dans $1+\mathbb{Q}_\ell[[t]]^+$].. Lorsque on admet les conjectures de Weil, et $F$ \inserted{est un motif,} $\underline{\chi}(F)$ se décompose suivant des poids en des $\sum(-1)^i\deleted{\unsure{\underline{\chi}}}\underline{h}^i_{\unsure{!}}(F)$, $\underline{h}^i(F)\in R(\mathbb{N},\mathbb{Q}_\ell)$, et la \deleted{\unsure{conna|connaissance}}connaissance de $L(X,F)(t)$ \uline{équivaut} à celle des $\underline{h}^i(F)$ : les $\underline{h}^i(F)$ correspondent à la décomposition de $L(X,F)$
% archive sheet number: 91
